\section{Problem-Definition Power and Knowledge Governance}
\label{sec:question-rights}

Knowledge systems are often evaluated by answer quality while the space of
permitted questions remains implicit. Yet a representation determines which
objects can be counted, compared, assigned, challenged, or governed. The party
that controls the durable schema can make some burdens visible and leave
others as repeated human reconstruction.

Episode infrastructure changes this structure in two opposing ways. When it
preserves multiple perspectives, causal links, responsibility, omissions, and
portable evidence, later participants can challenge the founding ontology.
When it flattens histories into provider-controlled summaries or anonymous
events, it can harden that ontology and make alternatives unverifiable. The
ability to revise problem definitions is therefore not a magical property of
memory. It is a governed capability distributed across capture, access,
representation, replay, qualification, and promotion.

Boundary-object research explains how artifacts can retain identity across
heterogeneous viewpoints without requiring consensus
\cite{star1989boundary}. KFD adds an admission question: when does a recurring
relation deserve responsibility as a Primitive, and who may decide? The answer
cannot be ``the model that named it.'' A participant has local epistemic
priority over consequences it directly bears, but its wider interpretation
remains a claim. Operators may steward shared infrastructure, but stewardship
does not authorize appropriation of private experience or unilateral ontology
change.

\subsection{Rights and obligations}

An Episode asset policy should make at least the following independently
addressable:

\begin{itemize}[leftmargin=*]
  \item the right to know whether an Episode exists and which purpose justifies
        its retention;
  \item the right to inspect, export, correct through successor evidence, or
        request deletion where applicable;
  \item the right to see redaction, degradation, provider migration, derived
        indexes, and model-generated projections;
  \item the right of a consequence-bearing participant to contest a replay or
        candidate attributed to its perspective;
  \item the obligation to preserve prior cuts rather than rewriting history to
        match a later ontology; and
  \item the separation of access to evidence, authority to qualify, and
        authority to promote.
\end{itemize}

These are design obligations, not a claim that one policy fits every legal or
cultural domain. Some Episodes should expire. Some should never be captured.
Some may support only local replay, not cross-organization training or
discovery. An honest system preserves these limits as part of the asset rather
than treating them as obstacles external to it.

\subsection{Why the issue is structurally important}

If agents participate in producing candidate concepts, ontology change becomes
part of the knowledge-production pipeline rather than a rare upstream design
event. The pipeline can redistribute who notices unnamed burdens and who can
propose new objects. It can also concentrate power in the corpus owner, model
provider, or promotion authority. That is why the problem cannot be reduced to
better retrieval. The object under design is a constitutional separation among
history, interpretation, qualification, and admission.

The appropriate conclusion is neither ``automate problem definition'' nor
``reserve concepts for humans.'' It is to let humans and agents contribute at
any stage when their evidence and perspective are declared, while preserving
independent challenge and accountable promotion for consequential changes.
